\documentclass[11pt]{article}

\usepackage[margin=0.85in]{geometry}
\usepackage[T1]{fontenc}
\usepackage[utf8]{inputenc}
\usepackage{amsmath}
\usepackage{booktabs}

\setlength{\parskip}{0.25em}
\setlength{\parindent}{0pt}

\title{Payload Bits per Frame}
\author{}
\date{11 August 2026}

\begin{document}
\maketitle

\section{Configuration}

An orthogonal frequency-division multiplexing (OFDM) frame carries data,
pilot symbols, a cyclic prefix, a synchronization signal, and a cyclic
redundancy check (CRC).  We calculate the payload for the current JUNA frame
at durations of one and two seconds, with a sampling rate of 9600 samples per
second.

The frame uses quadrature phase shift keying (QPSK), code rate 0.25, outer
and inner pilot spacings 5/5, and a cyclic prefix of 64 samples.  The
low-density parity-check (LDPC) code has check degree 14, the code horizon
fills the frame, and the receiver uses four partial-FFT parts.  The
synchronization signal uses 4096 samples, and the CRC uses 16 bits once per
frame.

\section{One-Second Frame}

Let $N$ be the fast Fourier transform size and let $B(N)$ be the number of
complete OFDM blocks in one frame.  The number of complete blocks is given by
\begin{equation}
  B(N)
  =
  \left\lfloor\frac{9600-4096}{N+64}\right\rfloor,
  \label{eq:blocks}
\end{equation}
where 9600 is the one-second sample budget, 4096 is the synchronization
overhead, and 64 is the cyclic prefix length.  Only complete blocks are
transmitted, so the samples remaining after the last complete block carry no
additional payload.

Let $k(N)$ be the LDPC information bits in one OFDM block after the outer
pilots and code rate are applied.  The user payload $P(N)$ is given by
\begin{equation}
  P(N)
  =
  B(N)k(N)
  -\left\lceil\frac{B(N)k(N)}{5}\right\rceil
  -16,
  \label{eq:payload}
\end{equation}
where the ceiling term gives the inner pilot bits and 16 is the frame CRC.
The cyclic prefix reduces the number of blocks that fit in the frame, while
the CRC is removed once from the complete frame payload.

Table~\ref{tab:payload} gives the payload returned by the current modem for
each FFT size.  The frame sample count includes the synchronization signal
and the cyclic prefix of every complete OFDM block.

\begin{table}[ht]
\centering
\caption{Payload bits per one-second frame after pilot, coding, cyclic prefix,
synchronization, and CRC overhead.}
\label{tab:payload}
\begin{tabular}{@{}rrrrrr@{}}
\toprule
$N$ & Blocks & Frame samples & LDPC $n/k$ & Before CRC & Payload \\
\midrule
 512 & 9 & 9280 & $816/204$   & 1468 & 1452 \\
1024 & 5 & 9536 & $1632/408$  & 1632 & \textbf{1616} \\
1536 & 3 & 8896 & $2448/612$  & 1468 & 1452 \\
2048 & 2 & 8320 & $3264/816$  & 1305 & 1289 \\
4096 & 1 & 8256 & $6544/1636$ & 1308 & 1292 \\
\bottomrule
\end{tabular}
\end{table}

The $N=1024$ configuration carries 1616 payload bits, which is the largest
payload among these five configurations.  The $N=512$ and $N=1536$
configurations both carry 1452 bits because their complete blocks contain the
same total number of LDPC information bits.

A larger FFT does not necessarily increase the payload in a fixed one-second
frame.  The synchronization signal, cyclic prefixes, and requirement for
complete OFDM blocks leave a different number of unused samples for each FFT
size.  These values describe transmitted frame capacity and do not measure
bit error rate or receiver success.

\section{Two-Second Frame}

The two-second frame has a budget of 19200 samples at the same sampling rate.
The synchronization signal and CRC are included once, so more of this budget
is available for complete OFDM blocks.

Let $B_2(N)$ be the number of complete OFDM blocks in the two-second frame.
The number of complete blocks is given by
\begin{equation}
  B_2(N)
  =
  \left\lfloor\frac{19200-4096}{N+64}\right\rfloor,
  \label{eq:blocks-two-seconds}
\end{equation}
where $N$ is the FFT size, 19200 is the two-second sample budget, 4096 is the
synchronization overhead, and 64 is the cyclic prefix length.  The larger
frame admits more complete OFDM blocks because synchronization is included
once.

Let $P_2(N)$ be the user payload in the two-second frame.  The payload is
given by
\begin{equation}
  P_2(N)
  =
  B_2(N)k(N)
  -\left\lceil\frac{B_2(N)k(N)}{5}\right\rceil
  -16,
  \label{eq:payload-two-seconds}
\end{equation}
where $k(N)$ is the LDPC information bits in one block, the ceiling term gives
the inner pilot bits, and 16 is the frame CRC.  This calculation removes the
inner pilots and CRC after the complete block count is fixed.

Table~\ref{tab:payload-two-seconds} gives the two-second payload returned by
the current modem for each FFT size.

\begin{table}[ht]
\centering
\caption{Payload bits per two-second frame after pilot, coding, cyclic prefix,
synchronization, and CRC overhead.}
\label{tab:payload-two-seconds}
\begin{tabular}{@{}rrrrr@{}}
\toprule
$N$ & Blocks & Frame samples & Before CRC & Payload \\
\midrule
 512 & 26 & 19072 & 4243 & 4227 \\
1024 & 13 & 18240 & 4243 & 4227 \\
1536 &  9 & 18496 & 4406 & 4390 \\
2048 &  7 & 18880 & 4569 & \textbf{4553} \\
4096 &  3 & 16576 & 3926 & 3910 \\
\bottomrule
\end{tabular}
\end{table}

The $N=2048$ configuration carries 4553 payload bits, which is the largest
payload among the five two-second configurations.  Thus the FFT size with the
largest payload changes when the frame duration changes.

At $N=1024$, two separate one-second frames carry $2\times1616=3232$ payload
bits.  One two-second frame carries 4227 bits, giving 995 additional bits
because synchronization and CRC are included once and three additional
complete OFDM blocks fit in the frame.

\section{Verification}

The current JUNA modem returned a valid configuration for all five FFT sizes
at both frame durations.  Retained result manifests independently record 1452 bits for
$N=512$, 1616 bits for $N=1024$, and 1289 bits for $N=2048$.  The $N=1536$
and $N=4096$ values, together with every two-second value, were obtained from
the same current modem calculation.

\end{document}
